\documentclass[10pt, letterpaper]{article}

\usepackage[
    top=1.2cm, bottom=0.75cm, left=2cm, right=2cm,
    footskip=1cm
]{geometry}
\usepackage[utf8]{inputenc}
\usepackage[T1]{fontenc}
\usepackage[spanish]{babel}
\usepackage{enumitem}
\usepackage{paracol}
\usepackage{titlesec}
\usepackage{hyperref}
\usepackage{xcolor}
\usepackage{charter}
\pagestyle{empty}

\definecolor{primaryColor}{RGB}{0,0,0}
\hypersetup{
    colorlinks=true,
    urlcolor=primaryColor
}

\titleformat{\section}{\large\bfseries}{}{0pt}{}[\titlerule]
\titlespacing{\section}{0pt}{0.55em}{0.4em}
\setlength{\columnsep}{1cm}
\setlist[itemize]{leftmargin=*, itemsep=2pt, topsep=2pt}
\setlength{\parindent}{0pt}

\begin{document}

%=====================
% ENCABEZADO
%=====================
\begin{center}
    {\LARGE \textbf{Agustina Loyola}} \\
    Data \& Operations | Fintech \& Crypto | Actuaria \\
    Buenos Aires, Argentina \\
    \href{mailto:agustinasol98@gmail.com}{agustinasol98@gmail.com} \quad | \quad +54 11 5989 0533 \\
    \end{center}
\vspace{0.1cm}

%=====================
% PERFIL
%=====================
\section*{Perfil Profesional}
Analista de datos y operaciones con formación actuarial (UBA, a una materia de graduarme), enfocada en fintech y crypto. Actualmente Operations Analyst en Roxom, en conciliación on-chain, reporting y analytics de producto. Trabajo con SQL (Athena/Trino, PostgreSQL), Python y R para análisis, automatización y construcción de dashboards y pipelines confiables.

%=====================
% EXPERIENCIA
%=====================
\section*{Experiencia Profesional}

\textbf{Operations Analyst} – Roxom \hfill \textit{Diciembre 2025 – Presente}
\begin{itemize}
    \item Conciliaciones diarias entre wallets y proveedores de la compañía (on-chain y exchanges).
    \item Reportes semanales de movimientos on-chain.
    \item Data analytics sobre los productos actuales de Roxom.
    \item Desarrollo de dashboards y queries en Metabase y SQL.
\end{itemize}
\vspace{0.1cm}

\textbf{ORT Argentina} \hfill \textit{Marzo 2023 – Presente} \\[3pt]
\textit{Docente de Economía Financiera} \hfill \textit{Marzo 2023 – Presente}
\begin{itemize}
    \item Docente de Macroeconomía y Finanzas en nivel secundario.
    \item Desarrollo de contenidos pedagógicos aplicados a casos reales.
\end{itemize}

\textit{Responsable del Programa de Educación Financiera} \hfill \textit{Marzo 2023 – Diciembre 2025}
\begin{itemize}
    \item Diseño e implementación del programa en nivel medio (alcance: +3000 estudiantes).
    \item Coordinación con equipo docente y oradores externos para articular contenidos y dinámica del programa.
\end{itemize}

\textit{Asistente de Dirección de Gestión de las Organizaciones} \hfill \textit{Marzo 2023 – Diciembre 2025}
\begin{itemize}
    \item Planificación e implementación de proyectos educativos en nivel medio.
    \item Coordinación entre equipos docentes y dirección para llevar adelante propuestas académicas.
\end{itemize}
\vspace{0.1cm}

\textbf{Finance Manager} – Tuareq (Industria Textil) \hfill \textit{Marzo 2023 – Agosto 2025}
\begin{itemize}
    \item Responsable del área financiera de una PyME textil: planificación, control y estrategia económica.
    \item Proyecciones anuales de ingresos, márgenes y crecimiento.
    \item Definición de estructura de costos, decisiones de inversión y análisis de rentabilidad.
    \item Reportes financieros ejecutivos para la dirección general.
\end{itemize}
\vspace{0.1cm}

\textbf{Co-Founder} – Esther Latam (Climate Tech) \hfill \textit{Mayo 2020 – Diciembre 2022}
\begin{itemize}
    \item Co-fundadora de una climate tech de reportes de impacto ambiental para organizaciones y eventos.
    \item Diseño del modelo de negocio y coordinación del desarrollo del MVP.
    \item Seguimiento financiero, presupuestos y reporting.
    \item Gestión de CRM, generación de pipeline comercial y alianzas estratégicas.
\end{itemize}

%=====================
% EDUCACIÓN
%=====================
\section*{Educación}

\textbf{Universidad de Buenos Aires (UBA)} \\
Actuaria en Economía · Estudiante avanzada (última materia para graduarme) \hfill \textit{2019 – Presente} \\
- Ayudante de cátedra en Estadística – Facultad de Ciencias Económicas

\vspace{0.1cm}

\textbf{ORT Argentina} \\
Bachiller Técnico en Administración de Empresas \hfill \textit{2012 – 2016}

%=====================
% HABILIDADES
%=====================
\section*{Habilidades Técnicas}
\begin{itemize}
    \item \textbf{Data \& Lenguajes:} SQL (Athena/Trino, PostgreSQL), Python (pandas, estadística), R, Excel avanzado
    \item \textbf{BI \& Reporting:} Metabase, dashboards, automatización de reportes
    \item \textbf{Finanzas:} forecasting, control presupuestario, métricas SaaS
    \item \textbf{Idiomas:} Inglés (C1)
\end{itemize}

\end{document}
